\pagestyle{empty}
\tight
\begin{tblr}{width=\textwidth,colspec={|Q[c,m,1.9cm]|X[l,m]|},hlines,vlines,hspan=minimal,row{1}={bg=headblue},rowsep=1pt,colsep=3pt}
\SetCell{c,m}\hfil\logo\hfil & \textbf{POLITEKNIK NEGERI MALANG}\par\textbf{JURUSAN TEKNOLOGI INFORMASI}\par\textbf{PROGRAM STUDI: D4 TEKNIK INFORMATIKA} \\
\SetCell[c=2]{c} \textbf{RENCANA PEMBELAJARAN SEMESTER (RPS)} & \\
\end{tblr}
\begin{longtblr}{width=\textwidth,colspec={|Q[l,m,1.9cm]|X[0.85,l,m]|X[1.45,l,m]|X[0.75,l,m]|X[1.15,l,m]|},hlines,vlines,hspan=minimal,rowsep=1pt,colsep=2pt,row{1,3}={bg=headgray,font=\bfseries}}
MATA KULIAH & KODE & BOBOT (SKS)/jam & SEMESTER & TGL. PENYUSUNAN \\
Bahasa Indonesia & RTI254003 & 2 SKS/4 jam & 4 & 1 Juli 2025 \\
OTORISASI & Dosen Pengembang RPS & Koordinator RMK & \SetCell[c=2]{l} Ka PRODI &  \\
 & Wilda Imama Sabilla, S.Kom., M.Kom. & Wilda Imama Sabilla, S.Kom., M.Kom. & \SetCell[c=2]{l}{\vspace{8mm}\par Dr. Ely Setyo Astuti, S.T., M.T.} &  \\
\SetCell[r=9]{l} Capaian Pembelajaran (CP) & \SetCell[c=4]{l,bg=headgray,font=\bfseries} Capaian Pembelajaran Lulusan Program Studi (CPL-Prodi) &  &  &  \\
 & CPL03 & \SetCell[c=3]{l} Mampu mengelola tim dan diri sendiri secara profesional serta mengomunikasikan dan mempresentasikan gagasan maupun hasil kerja secara efektif baik lisan maupun tulisan. &  &  \\
 & CPL04 & \SetCell[c=3]{l} Mampu menyusun deskripsi saintifik hasil kajian, pengembangan, atau implementasi teknologi informasi dan komunikasi dalam bentuk skripsi, laporan, atau artikel ilmiah. &  &  \\
 & \SetCell[c=4]{l,bg=headgray,font=\bfseries} Capaian Pembelajaran mata kuliah (CPL-MK) &  &  &  \\
 & CPMK0301 & \SetCell[c=3]{l} Mahasiswa mampu mengelola diri secara profesional, berkomunikasi efektif lisan dan tulisan dalam bahasa Indonesia baku untuk keperluan akademik dan profesional. &  &  \\
 & CPMK0401 & \SetCell[c=3]{l} Mahasiswa mampu menganalisis masalah dan menulis karya ilmiah berbasis data menggunakan kaidah bahasa Indonesia yang benar. &  &  \\
 & \SetCell[c=4]{l,bg=headgray,font=\bfseries} Kemampuan akhir tiap tahapan belajar (Sub-CPMK) &  &  &  \\
 & SCPMK0301-03001 & \SetCell[c=3]{l} Mahasiswa mampu menulis laporan teknis, makalah, dan proposal dalam bahasa Indonesia yang baku, sistematis, dan efektif. &  &  \\
 & SCPMK0401-03002 & \SetCell[c=3]{l} Mahasiswa mampu menyusun karya tulis ilmiah sederhana dengan rumusan masalah, metodologi, pembahasan, dan kesimpulan yang valid sesuai kaidah akademik. &  &  \\
\SetCell[r=2]{l} Penilaian & \SetCell[c=4]{l} *Nilai CPMK didapat dari agregasi nilai SUB-CPMK yang dibebankan &  &  &  \\
 & \SetCell[c=4]{l}{\tiny\begin{tblr}{width=\linewidth,colspec={|Q[c,m,0.1\linewidth]|Q[c,m,0.16\linewidth]|X[l,m]|X[l,m]|X[l,m]|X[l,m]|X[l,m]|X[l,m]|Q[c,m,0.08\linewidth]|},hlines,vlines,hspan=minimal,rowsep=1pt,colsep=0.7pt,row{1,2}={bg=headgray,font=\bfseries}}
Kode CPMK & Kode SCPMK & \SetCell[c=6]{c} Bobot per Bentuk Penilaian & & & & & & Total Bobot \\
 & & Tugas & Kuis 1 & UTS & Kuis 2 & PBL & UAS & \\
CPMK0301 & SCPMK0301-03001 & 15\%: laporan teknis dan proposal & 5\%: ejaan dan kalimat efektif & 8\%: laporan teknis & 0\% & 0\% & 10\%: presentasi dan karya tulis & 38\% \\
CPMK0401 & SCPMK0401-03002 & 15\%: karya tulis ilmiah & 0\% & 7\%: tata bahasa ilmiah & 5\%: penulisan ilmiah & 22\%: karya tulis ilmiah & 13\%: pengumpulan karya tulis & 62\% \\
\SetCell[c=8]{r}\textbf{Total Per Penilaian} & & & & & & & & \textbf{100\%} \\
\end{tblr}} &  &  &  \\
Diskripsi Singkat Mata Kuliah & \SetCell[c=4]{l} Bahasa Indonesia mengembangkan kemampuan berbahasa Indonesia baku untuk keperluan akademik dan profesional, meliputi penulisan laporan teknis, karya tulis ilmiah, dan presentasi di bidang TI. &  &  &  \\
Bahan Kajian / Materi Pembelajaran & \SetCell[c=4]{l}\begin{enumerate}\item Ragam bahasa, ejaan, diksi, dan kalimat efektif\item Paragraf, laporan teknis, proposal, dan presentasi ilmiah\item Karya tulis ilmiah: struktur, rumusan masalah, metodologi\item Sitasi, daftar pustaka, abstrak, dan artikel ilmiah\end{enumerate} &  &  &  \\
\SetCell[r=2]{l} Pustaka & \SetCell[c=4]{l} \textbf{Utama:}\begin{enumerate}\item Badan Pengembangan Bahasa. 2022. Pedoman Umum Ejaan Bahasa Indonesia (PUEBI Edisi V).\item Suparno \& Yunus, M. 2009. Keterampilan Dasar Menulis. UT.\item Dalman. 2018. Menulis Karya Ilmiah. Rajawali Press.\end{enumerate} &  &  &  \\
 & \SetCell[c=4]{l} \textbf{Pendukung:} Materi LMS, dokumentasi resmi, dan jobsheet praktikum. &  &  &  \\
Media Pembelajaran & \SetCell[c=2]{l} \textbf{Software:} LMS, pengolah kata, presentasi. &  & \SetCell[c=2]{l} \textbf{Hardware:} Komputer/laptop, proyektor. &  \\
Nama Dosen Pengampu & \SetCell[c=4]{l} Tim Pengajar Program Studi D4 TI &  &  &  \\
Matakuliah Syarat & \SetCell[c=4]{l} - &  &  &  \\
\end{longtblr}
\begin{longtblr}{width=\textwidth,colspec={|Q[c,m,0.06\textwidth]|X[1.35,l,m]|X[1.08,l,m]|X[1.16,l,m]|X[0.7,l,m]|X[1.24,l,m]|X[1.06,l,m]|X[0.98,l,m]|Q[c,m,0.05\textwidth]|},hlines,vlines,hspan=minimal,rowhead=2,row{1,2}={bg=headgreen,font=\bfseries},rowsep=1pt,colsep=1.5pt}
Mgg.\par Ke & Kemampuan Akhir Yang Direncanakan (Sub-CP-MK) & Bahan kajian (Materi Pembelajaran) & Bentuk dan Metode Pembelajaran & Estimasi Waktu & Pengalaman Belajar Mahasiswa & Kriteria \& Bentuk Penilaian & Indikator Penilaian & Bobot Penilaian (\%) \\
(1) & (2) & (3) & (4) & (5) & (6) & (7) & (8) & (9) \\
1 & SCPMK0301-03001 & Ragam bahasa: bahasa baku dan tidak baku. & Modalitas: Blended Learning. Bentuk: Luring/Daring. Metode: Case Method, diskusi, presentasi, dan peer review. & 1 x 2 x 50' tatap muka; 1 x 2 x 50' tugas/praktik mandiri. & Mahasiswa mengerjakan aktivitas terarah untuk ragam bahasa: bahasa baku dan tidak baku dan menunjukkan evidence hasil belajar. & Bentuk penilaian: Ragam Bahasa: Bahasa Baku dan Tidak Baku. Mengacu rubrik RTM. & Ketepatan bahasa; kualitas karya tulis; validasi dan dokumentasi. & 3 \\
2 & SCPMK0301-03001 & Ejaan yang Disempurnakan (EYD V) dan tanda baca. & Modalitas: Blended Learning. Bentuk: Luring/Daring. Metode: Case Method, diskusi, presentasi, dan peer review. & 1 x 2 x 50' tatap muka; 1 x 2 x 50' tugas/praktik mandiri. & Mahasiswa mengerjakan aktivitas terarah untuk ejaan yang disempurnakan (EYD V) dan tanda baca dan menunjukkan evidence hasil belajar. & Bentuk penilaian: Ejaan Yang Disempurnakan (EYD V) dan Tanda Baca. Mengacu rubrik RTM. & Ketepatan bahasa; kualitas karya tulis; validasi dan dokumentasi. & 3 \\
3 & SCPMK0301-03001 & Diksi dan kalimat efektif. & Modalitas: Blended Learning. Bentuk: Luring/Daring. Metode: Case Method, diskusi, presentasi, dan peer review. & 1 x 2 x 50' tatap muka; 1 x 2 x 50' tugas/praktik mandiri. & Mahasiswa mengerjakan aktivitas terarah untuk diksi dan kalimat efektif dan menunjukkan evidence hasil belajar. & Bentuk penilaian: Diksi dan Kalimat Efektif. Mengacu rubrik RTM. & Ketepatan bahasa; kualitas karya tulis; validasi dan dokumentasi. & 3 \\
4 & SCPMK0301-03001 & Kuis 1: ejaan dan kalimat efektif. & Modalitas: Blended Learning. Bentuk: Luring/Daring. Metode: Case Method, diskusi, presentasi, dan peer review. & 1 x 2 x 50' tatap muka; 1 x 2 x 50' tugas/praktik mandiri. & Mahasiswa mengerjakan aktivitas terarah untuk kuis 1: ejaan dan kalimat efektif dan menunjukkan evidence hasil belajar. & Bentuk penilaian: Kuis 1: Ejaan dan Kalimat Efektif. Mengacu rubrik RTM. & Ketepatan bahasa; kualitas karya tulis; validasi dan dokumentasi. & 5 \\
5 & SCPMK0301-03001 & Paragraf: kohesi dan koherensi. & Modalitas: Blended Learning. Bentuk: Luring/Daring. Metode: Case Method, diskusi, presentasi, dan peer review. & 1 x 2 x 50' tatap muka; 1 x 2 x 50' tugas/praktik mandiri. & Mahasiswa mengerjakan aktivitas terarah untuk paragraf: kohesi dan koherensi dan menunjukkan evidence hasil belajar. & Bentuk penilaian: Paragraf: Kohesi dan Koherensi. Mengacu rubrik RTM. & Ketepatan bahasa; kualitas karya tulis; validasi dan dokumentasi. & 3 \\
6 & SCPMK0301-03001 & Penulisan laporan teknis dan proposal. & Modalitas: Blended Learning. Bentuk: Luring/Daring. Metode: Case Method, diskusi, presentasi, dan peer review. & 1 x 2 x 50' tatap muka; 1 x 2 x 50' tugas/praktik mandiri. & Mahasiswa mengerjakan aktivitas terarah untuk penulisan laporan teknis dan proposal dan menunjukkan evidence hasil belajar. & Bentuk penilaian: Penulisan Laporan Teknis dan Proposal. Mengacu rubrik RTM. & Ketepatan bahasa; kualitas karya tulis; validasi dan dokumentasi. & 3 \\
7 & SCPMK0301-03001 & Presentasi dan diskusi ilmiah. & Modalitas: Blended Learning. Bentuk: Luring/Daring. Metode: Case Method, diskusi, presentasi, dan peer review. & 1 x 2 x 50' tatap muka; 1 x 2 x 50' tugas/praktik mandiri. & Mahasiswa mengerjakan aktivitas terarah untuk presentasi dan diskusi ilmiah dan menunjukkan evidence hasil belajar. & Bentuk penilaian: Presentasi dan Diskusi Ilmiah. Mengacu rubrik RTM. & Ketepatan bahasa; kualitas karya tulis; validasi dan dokumentasi. & 3 \\
8 & SCPMK0301-03001, SCPMK0401-03002 & UTS: laporan teknis dan tata bahasa. & Modalitas: Blended Learning. Bentuk: Luring/Daring. Metode: Case Method, diskusi, presentasi, dan peer review. & 1 x 2 x 50' tatap muka; 1 x 2 x 50' tugas/praktik mandiri. & Mahasiswa mengerjakan aktivitas terarah untuk UTS: laporan teknis dan tata bahasa dan menunjukkan evidence hasil belajar. & Bentuk penilaian: UTS: Laporan Teknis dan Tata Bahasa. Mengacu rubrik RTM. & Ketepatan bahasa; kualitas karya tulis; validasi dan dokumentasi. & 15 \\
9 & SCPMK0401-03002 & Karya tulis ilmiah: struktur dan jenis. & Modalitas: Blended Learning. Bentuk: Luring/Daring. Metode: Case Method, diskusi, presentasi, dan peer review. & 1 x 2 x 50' tatap muka; 1 x 2 x 50' tugas/praktik mandiri. & Mahasiswa mengerjakan aktivitas terarah untuk karya tulis ilmiah: struktur dan jenis dan menunjukkan evidence hasil belajar. & Bentuk penilaian: Karya Tulis Ilmiah: Struktur dan Jenis. Mengacu rubrik RTM. & Ketepatan bahasa; kualitas karya tulis; validasi dan dokumentasi. & 3 \\
10 & SCPMK0401-03002 & Rumusan masalah dan tinjauan pustaka. & Modalitas: Blended Learning. Bentuk: Luring/Daring. Metode: Case Method, diskusi, presentasi, dan peer review. & 1 x 2 x 50' tatap muka; 1 x 2 x 50' tugas/praktik mandiri. & Mahasiswa mengerjakan aktivitas terarah untuk rumusan masalah dan tinjauan pustaka dan menunjukkan evidence hasil belajar. & Bentuk penilaian: Rumusan Masalah dan Tinjauan Pustaka. Mengacu rubrik RTM. & Ketepatan bahasa; kualitas karya tulis; validasi dan dokumentasi. & 3 \\
11 & SCPMK0401-03002 & Metodologi penulisan ilmiah. & Modalitas: Blended Learning. Bentuk: Luring/Daring. Metode: Case Method, diskusi, presentasi, dan peer review. & 1 x 2 x 50' tatap muka; 1 x 2 x 50' tugas/praktik mandiri. & Mahasiswa mengerjakan aktivitas terarah untuk metodologi penulisan ilmiah dan menunjukkan evidence hasil belajar. & Bentuk penilaian: Metodologi Penulisan Ilmiah. Mengacu rubrik RTM. & Ketepatan bahasa; kualitas karya tulis; validasi dan dokumentasi. & 3 \\
12 & SCPMK0401-03002 & Kuis 2: penulisan ilmiah. & Modalitas: Blended Learning. Bentuk: Luring/Daring. Metode: Case Method, diskusi, presentasi, dan peer review. & 1 x 2 x 50' tatap muka; 1 x 2 x 50' tugas/praktik mandiri. & Mahasiswa mengerjakan aktivitas terarah untuk kuis 2: penulisan ilmiah dan menunjukkan evidence hasil belajar. & Bentuk penilaian: Kuis 2: Penulisan Ilmiah. Mengacu rubrik RTM. & Ketepatan bahasa; kualitas karya tulis; validasi dan dokumentasi. & 5 \\
13 & SCPMK0401-03002 & Sitasi dan daftar pustaka (APA/IEEE style). & Modalitas: Blended Learning. Bentuk: Luring/Daring. Metode: Case Method, diskusi, presentasi, dan peer review. & 1 x 2 x 50' tatap muka; 1 x 2 x 50' tugas/praktik mandiri. & Mahasiswa mengerjakan aktivitas terarah untuk sitasi dan daftar pustaka (APA/IEEE style) dan menunjukkan evidence hasil belajar. & Bentuk penilaian: Sitasi dan Daftar Pustaka (APA/IEEE Style). Mengacu rubrik RTM. & Ketepatan bahasa; kualitas karya tulis; validasi dan dokumentasi. & 3 \\
14 & SCPMK0401-03002 & Abstrak dan artikel ilmiah. & Modalitas: Blended Learning. Bentuk: Luring/Daring. Metode: Case Method, diskusi, presentasi, dan peer review. & 1 x 2 x 50' tatap muka; 1 x 2 x 50' tugas/praktik mandiri. & Mahasiswa mengerjakan aktivitas terarah untuk abstrak dan artikel ilmiah dan menunjukkan evidence hasil belajar. & Bentuk penilaian: Abstrak dan Artikel Ilmiah. Mengacu rubrik RTM. & Ketepatan bahasa; kualitas karya tulis; validasi dan dokumentasi. & 3 \\
15 & SCPMK0301-03001, SCPMK0401-03002 & Presentasi karya tulis ilmiah. & Modalitas: Blended Learning. Bentuk: Luring/Daring. Metode: Case Method, diskusi, presentasi, dan peer review. & 1 x 2 x 50' tatap muka; 1 x 2 x 50' tugas/praktik mandiri. & Mahasiswa mengerjakan aktivitas terarah untuk presentasi karya tulis ilmiah dan menunjukkan evidence hasil belajar. & Bentuk penilaian: Presentasi Karya Tulis Ilmiah. Mengacu rubrik RTM. & Ketepatan bahasa; kualitas karya tulis; validasi dan dokumentasi. & 3 \\
16 & SCPMK0401-03002 & PBL: penulisan karya tulis ilmiah. & Modalitas: Blended Learning. Bentuk: Luring/Daring. Metode: Case Method, diskusi, presentasi, dan peer review. & 1 x 2 x 50' tatap muka; 1 x 2 x 50' tugas/praktik mandiri. & Mahasiswa mengerjakan aktivitas terarah untuk PBL: penulisan karya tulis ilmiah dan menunjukkan evidence hasil belajar. & Bentuk penilaian: PBL: Penulisan Karya Tulis Ilmiah. Mengacu rubrik RTM. & Ketepatan bahasa; kualitas karya tulis; validasi dan dokumentasi. & 22 \\
17 & SCPMK0301-03001, SCPMK0401-03002 & UAS: presentasi dan pengumpulan karya tulis. & Modalitas: Blended Learning. Bentuk: Luring/Daring. Metode: Case Method, diskusi, presentasi, dan peer review. & 1 x 2 x 50' tatap muka; 1 x 2 x 50' tugas/praktik mandiri. & Mahasiswa mengerjakan aktivitas terarah untuk UAS: presentasi dan pengumpulan karya tulis dan menunjukkan evidence hasil belajar. & Bentuk penilaian: UAS: Presentasi dan Pengumpulan Karya Tulis. Mengacu rubrik RTM. & Ketepatan bahasa; kualitas karya tulis; validasi dan dokumentasi. & 23 \\
\end{longtblr}
